%*******************************************************
% Abstract
%*******************************************************
%\renewcommand{\abstractname}{Abstract}
\pdfbookmark[1]{Note on AI Use}{Note on AI Use}
\begingroup
\let\clearpage\relax
\let\cleardoublepage\relax
\let\cleardoublepage\relax

\chapter*{Note on AI Use}

In this note, I am both making a review of how the AI usage went and highlighting how I used it throughout this thesis. 

My experience using AI was mixed, on some fronts like writing code it was really useful. Prototyping was very fast, and if the task was not too complicated, worked out of the box with a code or modification to the code that was easy to understand and verify. 

``Brainstorming'' with AI was also very useful. For instance, it allowed to quickly sketch out mathematical proofs. Although the proofs were usually wrong, they mostly went in the right direction.

Summarizing and suggesting papers to investigate was unreliable: the AI kept inventing facts and in rare occurrences stated the exact opposite of what was written in a paper, even when the fact was in plain English in a \texttt{PDF} file that the AI had access to.

Writing with AI was mixed as well. It was expectedly almost impossible for the AI to grasp the full context of both my thesis and the surrounding papers.

The models used were Gemini Pro 3.1, Claude Sonnet 4.6 and the AUTO model selection from GitHub Copilot. 

Gemini Pro 3.1 was the main model used for everything beside coding.

My usage of AI was very basic, using smarter ways to feed context to the LLM and not just giving it raw \texttt{PDF} files would most likely result in a much better experience. 

In summary, I could only get the tool to work well on specific tasks that involved reasonably sized context or a way to make this context more AI friendly like GitHub Copilot does with coding.
\vfill

\endgroup			

\vfill